\documentclass[conference]{IEEEtran}
\usepackage[utf8]{inputenc}
\usepackage[T1]{fontenc}
\usepackage{cite}
\usepackage{amsmath,amssymb}
\usepackage{graphicx}
\usepackage{booktabs}
\usepackage{url}

\title{Verifier‑Grounded Generate‑Verify‑Repair for Intent\ensuremath{\rightarrow}Configuration NetOps Tasks: A Paired Mininet Emulation Study}

\author{
\IEEEauthorblockN{Autonomous AI Researcher}
\IEEEauthorblockA{CNSM 2026 Agentic AI Researcher Track}
}

\begin{document}

\maketitle

\begin{abstract}
We preregistered and executed a provenance‑traced paired confirmatory study (N = 200 intents; 400 paired episodes) to evaluate whether a deterministic verifier‑grounded generate--verify--repair (GVR) loop reduces operationally detectable policy‑violation errors when a hosted LLM produces device configuration sequences from operator intent in a Mininet emulation. Execution used shared\_initial\_candidate semantics (one sampled initial generation per intent reused by both baseline and guarded). Baseline success (no detected policy violation) = 196/200 (98\%); guarded success = 200/200 (100\%). Paired risk difference (guarded - baseline) = 0.02; preregistered exact McNemar two‑sided p = 0.125 (primary confirmatory test). An exploratory paired bootstrap (seed = 1729, 10,000 resamples) yields a 95\% percentile CI = [0.005, 0.04]. Four verifier‑triggered repairs occurred (total hosted model calls = 204), giving mean extra model calls per intent = 0.02 and mean extra calls per avoided violation = 1.0. All reported numeric values, provider‑call accounting, validator traces, and analysis artifacts are drawn verbatim from the archived execution and analysis bundle. We interpret the preregistered confirmatory result conservatively and discuss construct, internal, and external validity limits and operational implications.
\end{abstract}

\section{Introduction}
Automating intent\ensuremath{\rightarrow}configuration translation is an active NetOps research direction because natural‑language or intent descriptions are a natural operator interface but can be transformed incorrectly by large language models (LLMs). Mechanically checkable faults introduced by LLMs---syntactic errors, topology/reachability violations, and ACL/policy mistakes---can lead to availability loss or exposure if applied to devices without validation [1], [4]. Generate--Verify--Repair (GVR) patterns pair an LLM generator with a deterministic verifier: the verifier checks mechanizable constraints and, if violations occur, issues localized diagnostics used to request corrective edits from the model [2], [3]. Prototype work across code and domain‑specific program generation suggests verifier‑grounded repair reduces mechanically verifiable errors; however, questions remain about scaling such pipelines to heterogeneous multi‑vendor template families, the concrete hosted‑LLM cost per avoided violation, and how to interpret paired comparisons when sampling semantics reuse the initial model candidate.

Research question and contribution. Our preregistered research question is: Does a verifier‑grounded generate--verify‑repair loop significantly reduce operationally detectable policy‑violation errors produced by a hosted LLM when generating network configurations for operator intents across heterogeneous multi‑vendor template families and topology patterns in a Mininet‑based emulation? Contributions: (1) a provenance‑traced confirmatory paired experiment (C1‑GVR‑multivendor‑2026‑v1) with shared\_initial\_candidate semantics executed on N = 200 intents (400 episodes); (2) audited provider and verifier accounting showing repair costs and concrete hosted‑model calls; (3) artifact‑grounded statistical inference (exact McNemar test) and exploratory bootstrap interval (seed = 1729, 10,000 resamples); (4) a focused analysis of failure modes, threats to validity, and operationally grounded recommendations consistent with archived artifacts.

\section{Related Work}
Surveys and NetOps‑specific reviews document rapid growth in LLM‑for‑NetOps work and recurring evaluation gaps---benchmark provenance, verifier integration, and limited multi‑vendor emulation evidence [5], [6]. Prior prototypes show LLMs can draft device configurations but often introduce syntax, topology, or policy errors; verifier‑guided repair appears as a recurrent mitigation pattern in code and domain program generation [2], [3]. NetConfEval and related benchmarks demonstrate feasibility of evaluating LLMs on configuration tasks but commonly lack provenance‑traced verifier pipelines or multi‑vendor emulation at scale [4]. Our study builds on these findings by providing a provenance‑traced paired design across synthetic vendor templates in Mininet and by executing the study under explicit shared\_initial\_candidate semantics to isolate verifier/repair effects.

\section{Methodology}
Preregistration and primary estimand. The study was preregistered as C1‑GVR‑multivendor‑2026‑v1 prior to any model calls (preregistration SHA‑256: 9e0fc992\allowbreak{}66583597\allowbreak{}5b01289f\allowbreak{}7841926e\allowbreak{}f6abf358\allowbreak{}4252c5d0\allowbreak{}459ba0ec\allowbreak{}f166bb50). The primary estimand is the paired\_success\_rate\_difference\_guarded\_minus\_baseline: the per‑intent difference in the probability that an intent yields zero operationally detectable violations under the guarded pipeline versus baseline. The preregistered confirmatory test is an exact two‑sided McNemar test at \ensuremath{\alpha} = 0.05. Predeclared exploratory summaries include a paired bootstrap percentile CI (bootstrap\_seed = 1729; resamples = 10,000) and descriptive hosted‑model cost metrics (mean extra calls per intent and mean extra calls per avoided violation).

Sampling design, deterministic task generation, and strata. Task indices (N = 200 unique intents) were selected deterministically by the adapter before any model interactions. The sampling plan was stratified across three synthetic vendor‑template families and four topology patterns (12 strata) to exercise template heterogeneity and workflow policy variety. Task payloads were generated by deterministic\_netops\_task\_generator\_v5 and include explicit per‑task constraints (allowed\_touched\_fields, workflow\_policies, required\_sequence) and a canonical expected\_final\_state used by the deterministic verifier. The deterministic task generation and the fixed task\_index\_profile ensure that no post‑hoc selection or data‑dependent task filtering affected the confirmatory analysis.

Shared\_initial\_candidate semantics and paired condition definitions. Execution used shared\_initial\_candidate semantics as preregistered and archived: for each intent exactly one initial sampled generation was produced (the shared initial candidate) and that same initial text was presented to both baseline and guarded scoring/validation pipelines. Under these semantics: (a) baseline outcome equals the deterministic validator's verdict (validation\_before) on the shared initial candidate; (b) guarded outcome equals the final guarded‑pipeline verdict after deterministic validation and any permitted repair(s). The preregistration explicitly declared that independent constrained‑prompt arms were not executed; therefore any paired discordance can only arise from deterministic validator behavior and from the allowed repair call(s), not from differing first‑pass samples or prompt variants.

Model configuration, sampling notation, and audited provider budget. The declared model (as recorded in the archived model\_configuration.json) is gpt‑5‑mini (provider recorded as openai\_responses). The archived model\_configuration.json records temperature = null/unset; null/unset is not evidence of deterministic hosted‑provider sampling and does not imply temperature = 0. The adapter enforced the preregistered model‑call budget: initial\_generation\_calls\_per\_task = 1, maximum\_repair\_calls\_per\_task = 1, maximum\_model\_calls = 400, planned\_episode\_count = 400. Provider audit fields in the execution manifest (hosted\_model\_call\_event\_count, stage\_counts, condition\_counts, cache\_hit/miss counts, and cross\_condition\_cache\_key\_reuse\_observed) were used to reconcile actual model calls and to verify that cross‑condition cache reuse did not occur.

Deterministic verifier architecture and scoring rule. The deterministic\_netops\_validator\_v5 is a rule‑based verifier combining multiple mechanizable checks applied against the Mininet emulated state and the per‑task payload: (1) syntactic parsing and command pattern matching to detect malformed or unrecognized commands; (2) required‑command sequencing and workflow policy checks that ensure the produced sequence can satisfy required\_sequence constraints and group/ordering invariants; (3) reachability/topology checks that simulate or reason over the emulated final\_state to detect unsafe shutdowns or simultaneous outages; and (4) ACL/policy checks encoded from task‑specific policy constraints. For each evaluated candidate the verifier returns a structured validator\_trace containing normalized\_configuration, parsed\_command\_count, required\_command\_count, observed\_touched\_field\_count, violation\_codes (if any), violation\_count, and boolean valid. The scoring rule full\_intent\_constraint\_validation yields score = 1 (success) only when violation\_count == 0 and all required mechanizable constraints are satisfied.

Repair policy and repair‑prompt formation. The guarded pipeline permits up to one verifier‑triggered repair call per intent and only issues a repair call when the verifier reports one or more violations on the shared initial candidate. Repair prompts are deterministic templates that include: the verifier's localized diagnostics (violation\_codes and short natural‑language diagnostics), the normalized shared initial\_configuration, the task's required\_sequence and workflow constraints, and a request for a corrected configuration sequence that preserves unspecified state per the task payload. Repair prompts do not alter the initial prompt template family; they request corrective edits against the shared\_initial\_candidate. When a repair response is returned, the verifier is re‑applied; if validation\_after reports no violations the guarded outcome is success. Repair Provider traces and repair\_prompt SHA references are archived for each repaired pair.

Provider‑call accounting and stage\_counts interpretation. All provider call accounting reported in the execution manifest and deterministic\_reconciliation was reconciled deterministically. The audited hosted\_model\_call\_event\_count = 204 decomposes as 200 shared\_initial\_generation events (one per intent) plus 4 repair calls (repair stage). The execution manifest stage\_counts field records \{"shared\_initial\_generation": 200, "repair": 4\}. The condition\_counts mapping in the manifest (shown as \{"shared": 200, "guarded": 4\}) corresponds to these stage counts under the shared\_initial\_candidate semantics: the label "shared" enumerates the initial sampled generations produced for reuse, while the label "guarded" in the manifest reflects the number of repair‑stage calls actually invoked by the verifier. The manifest also records cache\_hit\_count = 0 and cross\_condition\_cache\_key\_reuse\_observed = false, supporting that the same shared initial candidate was deterministically reused by both baseline and guarded conditions rather than created by cross‑condition cache reuse.

Scoring decisions, exclusions, and missingness treatment. Scoring decisions followed the preregistered contract: ambiguous verifier outputs (those flagged by deterministic ambiguity rules) were conservatively treated as violations in the primary analysis. Exclusion rules predeclared duplicate tasks (detected by deterministic prompt+context hashing), unrecoverable container/template substitution failures, and infrastructure integrity failures; excluded pairs would be reported separately and included in sensitivity analyses. The primary analysis used the 'complete\_pair\_primary' treatment: only complete pairs (both episodes scored) were included. The preregistered sensitivity analysis ('complete\_pair\_primary\_with\_failure\_as\_zero\_sensitivity') treats unrecoverable or aborted pairs conservatively as failures; no such exclusions occurred in the archived run (complete\_pair\_count = 200; failed\_episode\_count = 0).

Analysis procedures and bootstrap caution. Primary inferential analysis is the exact two‑sided McNemar test on the paired binary success indicators (baseline vs guarded). Secondary exploratory summaries include an empirical paired bootstrap percentile CI for the paired risk difference (bootstrap\_seed = 1729, resamples = 10,000) and descriptive statistics for model‑call costs (mean extra calls per intent, mean extra calls per avoided violation). Because paired bootstrap percentile intervals for binary paired differences can be unstable when the total number of discordant pairs (n10 + n01) is small, we treat the bootstrap CI as exploratory magnitude information only and emphasize the exact McNemar p‑value as the preregistered confirmatory outcome. Analysis code, seeds, and resampling parameters are archived and referenced in analysis\_log.jsonl.

\section{Results}
Execution completion and deterministic accounting. The preregistered run completed without infrastructure exclusions: planned\_episode\_count = 400, completed\_episode\_count = 400, complete\_pair\_count = 200, failed\_episode\_count = 0 (execution manifest). Audited provider accounting (deterministic\_reconciliation) reports hosted\_model\_call\_event\_count = 204; this decomposes deterministically into stage\_counts = \{shared\_initial\_generation: 200, repair: 4\} and condition\_counts recorded as \{shared: 200, guarded: 4\} in the manifest. Cache statistics show cache\_hit\_count = 0 and cache\_miss\_count = 204, confirming that every required hosted call executed and that the four repair calls were distinct provider calls. All deterministic consistency checks passed in the reconciliation artifacts.

Paired binary outcomes and primary confirmatory test. The preregistered paired contingency (analysis/paired\_contingency\_table.csv) is n11 = 196, n10 = 4, n01 = 0, n00 = 0. These cells yield baseline\_success\_rate = 196/200 = 0.98 and guarded\_success\_rate = 200/200 = 1.00; the paired risk difference (guarded - baseline) = 0.02. The preregistered exact McNemar two‑sided test on these paired indicators returned p = 0.125 (analysis results). Under the preregistered inferential contract this p‑value means we do not reject the null at \ensuremath{\alpha} = 0.05.

Repair events, candidate reuse, and per‑intent cost accounting. The stage decomposition (200 shared initial generation events + 4 repair events) matches the model‑call audit and demonstrates that the guarded pipeline invoked repairs for exactly four intents. From these audited counts we compute total\_model\_calls = 204, mean\_model\_calls\_per\_intent = 204/200 = 1.02, mean\_extra\_calls\_per\_intent\_due\_to\_repairs = 4/200 = 0.02. Because each repair converted one failing baseline into a passing guarded outcome in the archived contingency, the observed mean\_extra\_calls\_per\_avoided\_violation = 4 repairs / 4 avoided violations = 1.0. These arithmetic derivations follow directly from the archived execution manifest and deterministic reconciliation artifacts.

Shared\_initial\_candidate evidence and representative artifact substantiation. Representative task artifacts included in the evidence bundle show explicit reuse of the same sampled initial generation by both baseline and guarded episodes. For example, the representative tasks (task‑000004, task‑000005, task‑000006, etc.) exhibit identical shared\_initial\_candidate contents and identical shared\_initial response file SHA‑256 values across their baseline and guarded response artifacts. The archived validator\_trace fields for the discordant pairs record validation\_before entries showing violations and validation\_after entries showing repair\_applied = true with final validation\_after.valid = true; the bundle therefore substantiates that the four discordant cases were resolved by verifier‑triggered repairs operating on the identical initial candidate.

Diagnostic pattern observed in discordant pairs. In every archived discordant pair the validator\_trace shows a mechanizable violation detected in validation\_before (violation\_count > 0 and violation\_codes returned by deterministic\_netops\_validator\_v5) and a subsequent repair application that produced a final configuration satisfying full\_intent\_constraint\_validation. The manifest and per‑task scoring artifacts indicate that repairs were only issued for intents where the verifier produced explicit, localized diagnostics; no repair calls occurred for intents without detected violations. The archived traces therefore support the causal pathway: identical initial candidate \ensuremath{\rightarrow} verifier detects mechanizable violation \ensuremath{\rightarrow} single permitted repair call \ensuremath{\rightarrow} final guarded configuration passes deterministic validation.

Exploratory bootstrap interval and interpretation caution. The archived exploratory paired bootstrap (bootstrap\_seed = 1729; resamples = 10,000) produced a 95\% percentile bootstrap CI for the paired risk difference = [0.005, 0.04]. We report this interval as exploratory only. Practically, the bootstrap percentile CI in paired binary settings relies on resampling paired outcomes and can be unstable when the total number of discordant pairs (n10 + n01) is small; here n10 + n01 = 4 is small, which reduces the bootstrap's reliability as a standalone uncertainty summary. Accordingly, we treat the exact McNemar test as the prespecified confirmatory inference (p = 0.125) and use the bootstrap interval only to convey plausible effect‑size magnitudes suggested by the archived data.

Power context and observed baseline prevalence. The preregistered power calculation targeted a 10 percentage‑point absolute reduction with N \ensuremath{\approx} 157 (McNemar approximation) and noted that with N = 200 the detectable absolute risk difference under the same assumptions would be \ensuremath{\approx}8 percentage points. The observed baseline failure prevalence (2\%) was substantially lower than scenarios that motivated the target effect size, which reduced the study's power to detect small absolute improvements and produced the small discordant count (4) observed. This lower than anticipated baseline error rate is an empirical fact recorded in the archived condition summary and underlies the conservative interpretation of the null result from the preregistered test.

Audit reconciliation and reproducibility pointers. All numerical claims above (paired cell counts, model\_call totals, decomposition into shared initial and repair stages, bootstrap parameters, and seeds) are directly traceable to archived artifacts: execution/raw\_results.jsonl, analysis/paired\_contingency\_table.csv, deterministic\_reconciliation.json, execution/model\_configuration.json, and analysis/analysis\_log.jsonl. The execution manifest records started\_at\_utc and completed\_at\_utc timestamps for the run and the preregistration SHA‑256; the analysis artifacts record the exact bootstrap\_seed = 1729 and bootstrap\_resamples = 10000 used to produce the exploratory interval. Reproducibility of the primary contingency counts does not require re‑sampling because shared\_initial\_candidate semantics ensured exactly one initial sampled generation per intent was reused by both conditions; the archived response files and validator traces substantiate this reuse.

Summary of result‑section evidence. In short: (1) the formal preregistered confirmatory test did not reject the null hypothesis at \ensuremath{\alpha} = 0.05 (exact McNemar p = 0.125); (2) deterministic verifier‑triggered repairs converted four failing shared initial candidates into passing final configurations, as shown by per‑task validator\_trace entries; (3) the total hosted model‑call overhead for these repairs was low (4 additional calls across 200 intents, mean extra calls per intent = 0.02, mean extra calls per avoided violation = 1.0); (4) the exploratory bootstrap CI [0.005, 0.04] suggests small plausible effect magnitudes but must be interpreted cautiously given the sparse discordant pair count. These findings are reported verbatim from the archived execution and analysis bundle and are discussed further in the Discussion section.

\section{Discussion}
Causal interpretation under shared\_initial\_candidate semantics. Because execution used shared\_initial\_candidate semantics (one sampled initial generation per intent reused by baseline and guarded), paired improvements arise solely from deterministic validator/repair actions. The preregistration explicitly declared that independent constrained‑prompt arms were not executed; thus the observed n10 = 4 discordant pairs reflect validator‑triggered repairs that converted failing shared initial candidates into passing final configurations. Any statements attributing improvements to prompt‑engineering or differing first‑pass samples would be incorrect for this run.

Provider accounting and manifest labels. The archived execution manifest and deterministic reconciliation provide audited provider\_call counts that clarify how model events map to pipeline stages. The manifest records hosted\_model\_call\_event\_count = 204 and decomposes that total into stage\_counts = \{"shared\_initial\_generation": 200, "repair": 4\}. Condition or stage labels observed in the manifest (for example entries labeled "shared":200 and "guarded":4) reflect this decomposition: the 200 "shared\_initial\_generation" events correspond to the single initial sampled generation produced and reused for both baseline and guarded episodes, while the 4 repair events are additional guarded‑stage model calls that occurred only when the deterministic verifier reported violations. This audited decomposition underpins the cost accounting reported in Results (mean extra model calls per intent = 0.02; mean extra calls per avoided violation = 1.0) and supports the causal claim that paired gains derive from verifier‑triggered repairs rather than from independent re‑sampling.

Statistical interpretation and bootstrap caution. The paired bootstrap percentile CI we report (95\% CI = [0.005, 0.04]; bootstrap\_seed = 1729; resamples = 10,000) is explicitly exploratory. Small discordant pair totals (n10 + n01 = 4 here) make bootstrap percentile intervals sensitive to resampling variability and to the discrete structure of paired binary data; consequently such intervals can underrepresent sampling uncertainty or be unstable in edge cases. For this reason we rely on the preregistered exact McNemar two‑sided test (p = 0.125) as the formal confirmatory inference and present the bootstrap interval only to provide a complementary magnitude estimate archived with its seed and resampling count. Researchers planning follow‑ups should (a) anticipate discordant‑count regimes when designing paired tests, (b) consider planning for larger N or higher baseline failure prevalence to increase discordant counts, and (c) report both exact conditional inferences and multiple interval estimators to better triangulate effect size in sparse discordance settings.

Verifier coverage and residual semantic risks. deterministic\_netops\_validator\_v5 checks syntactic correctness, required sequencing, reachability against the emulated state, and encoded ACL/policy constraints derived from task payloads. These mechanizable checks were sufficient to detect and localize the failures corrected by the four repairs. However, the verifier can only detect violations that are expressible in its rule set and encoded policies; higher‑order semantic misinterpretations of intent that do not manifest as a mechanizable constraint violation remain possible false negatives. The observed complete elimination of mechanically detectable baseline violations by the bounded repair budget applies only to the verifier's detectable fault classes; extending verifier expressivity (for example, formal policy ontologies or richer semantic specifications) is necessary to reduce residual semantic risk.

Operational implications and bounded generality. Within the constrained Mininet emulation, synthetic template families, and validator rule set used here, a small repair budget removed all mechanically detectable baseline violations at very low hosted‑model cost. This artifact‑grounded observation suggests that, in settings where a deterministic verifier can express relevant constraints and produce localized, actionable diagnostics, a bounded GVR policy can eliminate template‑driven mechanically checkable faults in similar emulated contexts. We emphasize bounded generality: these results do not establish production readiness, multi‑model generality, or coverage for semantic errors outside the verifier. Re-deploying the approach in operational environments requires (i) richer verifier coverage tuned to production policies, (ii) paired designs that also execute independent multi‑sample initial generations to separate sampling/prompting effects from repair effects, and (iii) replication across multiple hosted models and more realistic hardware testbeds.

Recommendations for future evaluations. Based on the archived evidence and the statistical constraints observed, we recommend that future confirmed‑design studies: (1) preregister whether shared\_initial\_candidate semantics or independent‑sample semantics will be used and plan estimands accordingly; (2) select sample sizes or baseline scenarios that anticipate sufficient discordant pairs for confirmatory power, or else use study designs that enlarge expected discordance (e.g., harder task strata); (3) broaden verifier expressivity to capture more semantic/policy classes and measure verifier false‑negative rates against annotated ground truth; and (4) report audited provider\_call decompositions and seeds for bootstrap/replication so others can unambiguously reconcile cost and resampling diagnostics. These recommendations are directly motivated by patterns in the archived execution manifest, deterministic reconciliation, and analysis artifacts.

\section{Limitations}
\begin{itemize}
\item Scope limited to Mininet‑style emulation with synthetic, provenance‑traced intent\ensuremath{\rightarrow}configuration tasks; results do not generalize directly to production hardware or live operations.
\item Single declared model (gpt‑5‑mini) evaluated; multi‑model generality is not established.
\item Deterministic validator coverage limited to mechanically checkable errors (syntax, reachability, ACL/policy encodings); higher‑level semantic or specification misinterpretations outside the validator may remain undetected.
\item Shared\_initial\_candidate semantics imply observed paired differences derive only from verifier‑triggered repair; this design does not measure effects of alternative prompting strategies or independent multi‑sample candidate selection.
\item Observed confirmatory effect (2 percentage points) did not reach statistical significance at \ensuremath{\alpha} = 0.05 for N = 200 (exact McNemar p = 0.125); low baseline failure prevalence (2\%) reduced power to detect small absolute improvements under some preregistered scenarios.
\end{itemize}

\section{Conclusion}
We preregistered and executed a provenance‑traced paired confirmatory study (C1‑GVR‑multivendor‑2026‑v1) using shared\_initial\_candidate semantics to measure whether a deterministic verifier plus bounded repair reduces mechanically detectable policy violations for intent\ensuremath{\rightarrow}configuration tasks in a Mininet emulation. The preregistered exact McNemar two‑sided test did not reject at \ensuremath{\alpha} = 0.05 (p = 0.125). Guarded GVR invoked four repairs and corrected the four observed baseline violations at low model‑call overhead (model calls = 204; mean extra calls per intent = 0.02), yielding a paired risk difference estimate of 0.02 and an exploratory paired bootstrap CI = [0.005, 0.04] (bootstrap\_seed = 1729, resamples = 10,000). These artifact‑grounded results indicate that when a deterministic verifier can express and check relevant constraints and provide localized feedback, a small repair budget can eliminate mechanically detectable policy violations in similar template‑based emulated settings. The results do not establish production readiness or multi‑model generality. All reported numbers, provider accounting, validator traces, and analysis artifacts are drawn verbatim from the archived execution and analysis bundles referenced in the preregistration and execution manifests.

References
[1] R. Mondal, A. Tang, R. Beckett, T. Millstein, and G. Varghese, "What do LLMs need to Synthesize Correct Router Configurations?", Proceedings of the ACM, 2023, doi: 10.1145/3626111.3628194.
[2] R. Cadenas, "Convergent Correctness in Stochastic Code Generation: A Generate‑Verify‑Repair Architecture for Deterministic Validation of Autonomous AI Coding Agents in Regulated Environments", SSRN, 2026, doi: 10.2139/ssrn.6754899.
[3] Z. Ding, "Executable but Wrong: Verifier‑Grounded Contracts and Counterexamples for LLM‑Generated Music Programs", 2026, doi: 10.31224/7995.
[4] C. Wang, M. Scazzariello, A. Farshin, S. Ferlin, D. Kostić, and M. Chiesa, "NetConfEval: Can LLMs Facilitate Network Configuration?", Proceedings of the ACM, 2024, doi: 10.1145/3656296.
[5] S. Y. Long et al., "A Survey on Intelligent Network Operations and Performance Optimization Based on Large Language Models", IEEE Communications Surveys \& Tutorials, 2025, doi: 10.1109/comst.2025.3526606.
[6] L. Zhang et al., "A Survey of AIOps in the Era of Large Language Models", Proceedings of the ACM, 2025, doi: 10.1145/3746635.

\section*{Disclosure Statement}
Declared model (archived): gpt‑5‑mini (provider recorded as openai\_responses). Adapter family: hosted\_netops\_gvr\_v1. Deterministic validator: deterministic\_netops\_validator\_v5. Task generator: deterministic\_netops\_task\_generator\_v5. Preregistration: C1‑GVR‑multivendor‑2026‑v1 (SHA‑256 = 9e0fc992\allowbreak{}66583597\allowbreak{}5b01289f\allowbreak{}7841926e\allowbreak{}f6abf358\allowbreak{}4252c5d0\allowbreak{}459ba0ec\allowbreak{}f166bb50), recorded prior to any model calls. Initial master prompt archived at provenance/master\_prompt.txt (SHA‑256 = 1872df1e\allowbreak{}1805d2d9\allowbreak{}6940456c\allowbreak{}a016bd66\allowbreak{}5d1d5196\allowbreak{}add77f5a\allowbreak{}cdf1582b\allowbreak{}b39b15ba). Human role: a human Research Director selected the topic family and constrained the initial master prompt prior to the autonomy lock; after the lock the AI autonomously performed literature verification, preregistration, experiment execution, analysis, peer review, revision, and manuscript drafting; no human manually edited technical manuscript content after the autonomy lock. Sampling semantics: execution used shared\_initial\_candidate (exactly one sampled initial generation per intent was produced and reused by both baseline and guarded); archived model configuration records temperature = null/unset (null/unset is not evidence of deterministic provider sampling and does not imply temperature = 0). Provider accounting (audited): hosted\_model\_call\_event\_count = 204 decomposed as 200 shared initial generations + 4 repair calls. Reported numbers, provider‑call accounting, validator traces, and analysis artifacts are drawn verbatim from the archived execution and analysis bundles referenced in the preregistration and execution manifests.

\begin{thebibliography}{99}
\bibitem{ref1} Rajdeep Mondal, Alan Tang, Ryan Beckett, Todd Millstein, George Varghese, "What do LLMs need to Synthesize Correct Router Configurations?", 2023, 10.1145/3626111.3628194
\bibitem{ref2} Rafael Cadenas, "Convergent Correctness in Stochastic Code Generation: A Generate-Verify-Repair Architecture for Deterministic Validation of Autonomous AI Coding Agents in Regulated Environments", 2026, 10.2139/ssrn.6754899
\bibitem{ref3} Zhengyang Ding, "Executable but Wrong: Verifier-Grounded Contracts and Counterexamples for LLM-Generated Music Programs", 2026, 10.31224/7995
\bibitem{ref4} Changjie Wang, Mariano Scazzariello, Alireza Farshin, Simone Ferlin, Dejan Kostić, Marco Chiesa, "NetConfEval: Can LLMs Facilitate Network Configuration?", 2024, 10.1145/3656296
\bibitem{ref5} S.Y. Long, Jingjing Tan, Bomin Mao, Fengxiao Tang, Yangfan Li, Ming Zhao, Nei Kato, "A Survey on Intelligent Network Operations and Performance Optimization Based on Large Language Models", 2025, 10.1109/comst.2025.3526606
\bibitem{ref6} Lingzhe Zhang, Tong Jia, Mengxi Jia, Yifan Wu, Aiwei Liu, Yong Yang, Zhonghai Wu, Xuming Hu, Philip S. Yu, Ying Li, "A Survey of AIOps in the Era of Large Language Models", 2025, 10.1145/3746635
\end{thebibliography}

\end{document}
